\chapter*{Abstract}
\addcontentsline{toc}{chapter}{Abstract}

The rapid globalization of the workforce has amplified the necessity for clear, standardized English communication. For non-native English speakers, particularly those in India, mother-tongue influence (MTI) heavily affects phonology, leading to communication barriers despite possessing strong grammatical and lexical knowledge. Traditional Computer-Assisted Pronunciation Training (CAPT) systems and generalized Automatic Speech Recognition (ASR) engines typically rely on acoustic models trained predominantly on Western accents (US/UK). As a result, these systems unfairly penalize natural phonetic characteristics inherent in Indian English, such as retroflex stops and dental consonants, leading to high error rates and user frustration.

This internship report presents the research, development, and deployment of \textbf{VoiceScore}, a highly specialized, AI-powered Indian English pronunciation scoring and diagnostic system. Unlike generic dictation ASRs, VoiceScore is engineered to act as a phoneme-level acoustic diagnostic coach. It leverages a Wav2Vec2 self-supervised backbone integrated with a Contrastive Connectionist Temporal Classification (CTC) head. This architecture projects acoustic frames and target phoneme embeddings into a shared 256-dimensional metric space, allowing the system to robustly measure acoustic distances and generate precise Goodness of Pronunciation (GoP) scores without relying on explicit frame pre-segmentation.

A critical contribution of this project was the development of a comprehensive, 133,737-word Indian English phonological lexicon. By applying context-sensitive sliding-window rules, standard American ARPAbet transcriptions were accurately mapped to a narrow 60-token International Phonetic Alphabet (IPA) inventory designed specifically for the Indian linguistic context. This process addressed intricate phenomena such as palatalization, labialization, and retroflexion.

The system was trained across two major sessions, evolving from a narrow academic domain dataset (NPTEL) to a robust, multi-accented mixture including Common Voice India, Svarah, and MUCS datasets. Technical challenges involving CP1252 double-encoding mojibake, data loader shared memory exhaustion (OOM), and C-level segmentation faults during forced alignment were systematically diagnosed and resolved. The final production model achieved a highly competitive Mean Phoneme Error Rate (PER) of 14.17\% on the validation split and demonstrated excellent generalization with a median PER of 16.28\% over an extensive 17,808-sample test suite, exhibiting a 0.00\% representation collapse rate.

Finally, the project culminated in the deployment of a modern, full-stack interactive web application. Deployed using a split-architecture design across Vercel (Next.js frontend) and Hugging Face Spaces (FastAPI backend with a Supabase PostgreSQL database), the product features real-time word heatmaps, interactive practice drills, and an AI-driven conversational tutor, providing users with actionable, stress-free pronunciation feedback.
